\chapter{Results}
\label{cha:results}

In this chapter, we present the results obtained from the fire risk classification 
experiments. The analysis is structured to evaluate the performance of individual 
models before moving to the ensemble approach. The primary goal is to assess the 
ability of different architectures to classify fire risk levels accurately across 
various input configurations and training strategies.

We evaluate three deep learning architectures individually:
\begin{itemize}
    \item ResNet50
    \item DenseNet161
    \item EfficientNet-B5
\end{itemize}

Each model is tested using one of the three different input configurations:
\begin{itemize}
    \item \gls{irg} images
    \item \gls{rgb} images
    \item \gls{rgb} images with \gls{ndvi} as a fourth channel
\end{itemize}

These configurations allow us to determine the most effective input format 
for fire risk classification. The models are trained across multiple \glspl{lr} and \glspl{wd} using 
the AdamW optimizer, which is widely used in 
deep learning due to its improved stability and convergence properties compared 
to traditional optimizers. Given its effectiveness, we did not experiment with 
alternative optimizers. If the results for an architecture are unsatisfactory, 
we further experiment with a gradual unfreezing strategy, as is the case for 
DenseNet161. It is also important to note that the models results are averaged
over three runs to ensure robustness and reliability.

Following the evaluation of individual architectures, we present the results of 
different ensemble models combining ResNet50, DenseNet161, and EfficientNet-B5 to 
leverage the strengths of each network. The ensemble approach is expected to enhance 
classification performance by integrating complementary decision-making processes 
from the individual models.

Finally, we analyze failure cases in the ensemble model's final results. Specifically, 
we focus on instances where all models in the ensemble incorrectly predict an average 
fire risk when, in reality, a very high fire risk is present. Understanding these 
failure cases is crucial for improving model robustness and minimizing false negatives 
in real-world applications.

\section{ResNet50 Results}

%                                                                  %
% ------------------------- RESNET50 \gls{irg} ------------------------- %
%                                                                  %
The initial experiments aim to evaluate how well ResNet50 adapts to the images across different 
\glspl{lr}. Figure \ref{fig:r50_results} presents the training and validation performance for 
\gls{rgb} images, showing trends in both loss and accuracy. Since the trends for \gls{irg} and \gls{rgb}+\gls{ndvi} 
images were similar, only the \gls{rgb} graphs are displayed for clarity.

\begin{figure}[h]
    \centering
    \includegraphics[width=1\textwidth]{../Figures/resnet50_averages.png}
    \caption{Average training and validation performance of ResNet50 across different \glspl{lr} with \gls{rgb} images only.}
    \label{fig:r50_results}
\end{figure}

As seen in Figure \ref{fig:r50_results}, models trained with higher \glspl{lr} 
show more unstable validation loss. In particular, $1\times10^{-4}$ shows significant 
fluctuations, indicating instability and potential overfitting. Despite achieving lower 
training loss, its validation accuracy stagnates, reinforcing the idea that overly 
aggressive \glspl{lr} lead to poor generalization.

In contrast, lower \glspl{lr} ($1\times10^{-5}$ and $5\times10^{-6}$) demonstrate 
smoother and more stable validation loss trends. This suggests that smaller 
\glspl{lr} enable a more controlled optimization process, reducing drastic 
fluctuations and improving generalization. Overall, the \glspl{lr} of $1\times10^{-5}$ 
and $5\times10^{-6}$ appear to have the best balance between validation and training performance.

\subsection{IRG Results}

The results presented in Table \ref{tab:resnet50_irg} provide clearer
insights into the model's performance, with the input images being \gls{irg} images.

\begin{table}[h]
    \centering
    \renewcommand{\arraystretch}{1.3}
    \setlength{\tabcolsep}{8pt}
    \begin{tabular}{|c|cccccc|}
        \hline
        & \multicolumn{2}{c}{\textbf{Validation}} & \multicolumn{2}{c}{\textbf{Train}} & \multicolumn{2}{c|}{\textbf{Test}} \\
        \textbf{\gls{lr}} & Accuracy & Loss & Accuracy & Loss & Accuracy & F1 Score \\
        \hline
        1e-4 & 75.86          & 0.4964 & 83.87 & 0.3552 & 75.39          & 0.6354 \\
        5e-5 & 75.84          & 0.4973 & 84.44 & 0.3583 & 75.90          & 0.7280 \\
        1e-5 & \textbf{83.62} & 0.3675 & 87.18 & 0.3011 & 80.13          & 0.7621 \\
        5e-6 & 79.45          & 0.4260 & 82.07 & 0.3944 & \textbf{80.82} & 0.7823 \\
        \hline
    \end{tabular}
    \caption{Average results for Resnet50 with \gls{irg} images across different \glspl{lr}.}
    \label{tab:resnet50_irg}
\end{table}

One key observation from Table \ref{tab:resnet50_irg} is that lower learning 
rates tend to yield better validation performance. A learning rate of 
$1\times10^{-5}$ achieves the highest validation accuracy (\textbf{83.62\%}), outperforming 
higher values such as $1\times10^{-4}$ (\textbf{75.86\%}) and $5\times10^{-5}$ (\textbf{79.45\%}). 
This reinforces the idea that larger \glspl{lr} may cause the model to converge too quickly to 
suboptimal solutions, limiting generalization.

Interestingly, the lowest learning rate ($5\times10^{-6}$) achieves the highest test 
accuracy (\textbf{80.82\%}), indicating that very small \glspl{lr} facilitate 
more stable convergence. However, its validation accuracy does not match that of 
$1\times10^{-5}$, which suggests a possible trade-off between validation and test performance.

The effect of \gls{wd} on the model's performance is then examined, with the
results summarized in Table \ref{tab:resnet50_wd_irg} for a learning rate of $5\times10^{-6}$.

\begin{table}[h]
    \centering
    \renewcommand{\arraystretch}{1.3}
    \setlength{\tabcolsep}{8pt}
    \begin{tabular}{|c|cccccc|}
        \hline
        & \multicolumn{2}{c}{\textbf{Validation}} & \multicolumn{2}{c}{\textbf{Train}} & \multicolumn{2}{c|}{\textbf{Test}} \\
        \textbf{\gls{wd}} & Accuracy & Loss & Accuracy & Loss & Accuracy & F1 Score \\
        \hline
        1e-2 & \textbf{81.32} & 0.4313 & 83.94 & 0.3568 & \textbf{82.90} & 0.8024 \\
        1e-4 & 80.17 & 0.4179 & 81.06 & 0.4106 & 80.31 & 0.7711 \\
        1e-6 & 79.88 & 0.4185 & 83.08 & 0.3918 & 80.57 & 0.7791 \\
        0    & 79.45 & 0.4260 & 82.07 & 0.3940 & 80.82 & 0.7823 \\
        \hline
    \end{tabular}
    \caption{Average results for Resnet50 with \gls{irg} images across different \gls{wd} values for a learning rate of 5e-6.}
    \label{tab:resnet50_wd_irg}
\end{table}

The \gls{wd} of $1\times10^{-2}$ 
achieves the highest test accuracy (\textbf{82.90\%}) and test F1 score 
(\textbf{0.8024}). This is a clear improvement over using no \gls{wd}, 
which results in a test accuracy of \textbf{80.82\%} and a test F1 score 
of \textbf{0.7820}. This suggests that applying \gls{wd} likely 
improves generalization by introducing regularization during training.

The lower \gls{wd} of $1\times10^{-6}$ achieves a similar test accuracy (\textbf{80.57\%}) 
but does not outperform $1\times10^{-2}$. This indicates that more regularization is beneficial, excessively small 
\gls{wd} values may not provide significant advantages. Overall, a \gls{wd} of $1\times10^{-2}$ appears to be 
the best choice for this setting.


%                                                                  %
% ------------------------- RESNET50 RGB ------------------------- %
%                                                                  %

\subsection{RGB Results}

The results for ResNet50 trained on \gls{rgb} images (Table \ref{tab:resnet50_rgb}) 
follow a similar trend to the \gls{irg} experiments, where lower \glspl{lr} 
generally yield better validation and test performance. The highest validation 
accuracy of \textbf{80.75\%} is achieved at a learning rate of $5\times10^{-6}$, with a 
corresponding test accuracy of \textbf{80.05\%} and an F1 score of \textbf{0.7736}. 

\begin{table}[h]
    \centering
    \renewcommand{\arraystretch}{1.3}
    \setlength{\tabcolsep}{8pt}
    \begin{tabular}{|c|cccccc|}
        \hline
        & \multicolumn{2}{c}{\textbf{Validation}} & \multicolumn{2}{c}{\textbf{Train}} & \multicolumn{2}{c|}{\textbf{Test}} \\
        \textbf{\gls{lr}} & Accuracy & Loss & Accuracy & Loss & Accuracy & F1 Score \\
        \hline
        1e-4 & 75.72 & 0.4787 & 77.90 & 0.4684 & 75.91 & 0.7436 \\
        5e-5 & 75.43 & 0.5077 & 68.57 & 0.5719 & 78.24 & 0.7469 \\
        1e-5 & 80.60 & 0.4246 & 80.63 & 0.4151 & 78.50 & 0.7562 \\
        5e-6 & 80.75 & 0.4131 & 81.07 & 0.3983 & 80.05 & 0.7736 \\
        \hline
    \end{tabular}
    \caption{Average results for Resnet50 with \gls{rgb} images across different \glspl{lr}.}
    \label{tab:resnet50_rgb}
\end{table}

This suggests once more that smaller \glspl{lr} contribute to more stable training and improved 
generalization. In contrast, higher \glspl{lr} such as $1\times10^{-4}$ and 
$5\times10^{-5}$ exhibit lower validation accuracy (\textbf{75.72\%} and \textbf{75.43\%}, respectively) 
and higher validation loss, indicating less stable optimization. Notably, the learning 
rate of $5\times10^{-6}$ provides a strong balance between validation accuracy (\textbf{80.75\%}) 
and test accuracy (\textbf{80.05\%}), similar to the best-performing \gls{irg} configuration.

The \gls{wd} experiments (Table \ref{tab:resnet50_wd_rgb}) show that moderate 
regularization improves performance, similar to the \gls{irg} results. A \gls{wd} of 
$1\times10^{-2}$ achieves the highest validation accuracy (\textbf{82.18\%}) and the highest test accuracy 
(\textbf{80.83\%}), reinforcing the importance of regularization for improved generalization. 

\begin{table}[h]
    \centering
    \renewcommand{\arraystretch}{1.3}
    \setlength{\tabcolsep}{8pt}
    \begin{tabular}{|c|cccccc|}
        \hline
        & \multicolumn{2}{c}{\textbf{Validation}} & \multicolumn{2}{c}{\textbf{Train}} & \multicolumn{2}{c|}{\textbf{Test}} \\
        \textbf{\gls{wd}} & Accuracy & Loss & Accuracy & Loss & Accuracy & F1 Score \\
        \hline
        1e-2 & 82.18 & 0.4173 & 82.36 & 0.3733 & 80.83 & 0.7730 \\
        1e-4 & 79.43 & 0.4255 & 80.41 & 0.3832 & 80.31 & 0.7680 \\
        1e-6 & 80.45 & 0.4482 & 82.79 & 0.3767 & 80.42 & 0.7710 \\
        0    & 80.75 & 0.4131 & 81.07 & 0.3983 & 80.05 & 0.7736 \\
        \hline
    \end{tabular}
    \caption{Average results for Resnet50 with \gls{rgb} images across different \gls{wd} values for a learning rate of 5e-6.}
    \label{tab:resnet50_wd_rgb}
\end{table}


The model trained with no \gls{wd} still performs relatively 
well, with a test accuracy of \textbf{80.05\%} and an F1 score of \textbf{0.7736}. However, as seen 
previously, \gls{wd} helps to control overfitting and stabilize validation performance. 


%                                                                         %
% ------------------------- RESNET50 RGB + NDVI ------------------------- %
%                                                                         %

\subsection{RGB + NDVI Results}

In this section, we analyze the impact of incorporating \gls{ndvi} alongside \gls{rgb} inputs 
on model performance.

\begin{table}[h]
    \centering
    \renewcommand{\arraystretch}{1.3}
    \setlength{\tabcolsep}{8pt}
    \begin{tabular}{|c|cccccc|}
        \hline
        & \multicolumn{2}{c}{\textbf{Validation}} & \multicolumn{2}{c}{\textbf{Train}} & \multicolumn{2}{c|}{\textbf{Test}} \\
        \textbf{\gls{lr}} & Accuracy & Loss & Accuracy & Loss & Accuracy & F1 Score \\
        \hline
        1e-4 & 77.73 & 0.4566 & 82.83 & 0.3851 & 74.35 & 0.6883 \\
        5e-5 & 79.17 & 0.4572 & 77.90 & 0.4603 & 72.80 & 0.7295 \\
        1e-5 & \textbf{80.17} & 0.4682 & 84.67 & 0.3448 & 77.47 & 0.6863 \\
        5e-6 & 78.59 & 0.4455 & 80.24 & 0.4151 & 77.72 & 0.6958 \\
        \hline
    \end{tabular}
    \caption{Average results for Resnet50 with \gls{rgb} images + \gls{ndvi} across different \glspl{lr}.}
    \label{tab:resnet50_rgb_ndvi}
\end{table}

When incorporating \gls{ndvi} alongside \gls{rgb} inputs, the results (Table \ref{tab:resnet50_rgb_ndvi}) 
show some performance variation across \glspl{lr}. The highest validation accuracy 
(\textbf{80.17\%}) is achieved at $1\times10^{-5}$, with a corresponding test accuracy of 
\textbf{77.47\%}.

The lowest learning rate ($5\times10^{-6}$) results in a test accuracy of \textbf{77.72\%} 
and an F1 score of \textbf{0.6958}. Additionally, higher \glspl{lr}, such as $1\times10^{-4}$, 
show slightly lower validation and test accuracies (\textbf{77.73\%} and \textbf{74.35\%}).

\begin{table}[h]
    \centering
    \renewcommand{\arraystretch}{1.3}
    \setlength{\tabcolsep}{8pt}
    \begin{tabular}{|c|cccccc|}
        \hline
        & \multicolumn{2}{c}{\textbf{Validation}} & \multicolumn{2}{c}{\textbf{Train}} & \multicolumn{2}{c|}{\textbf{Test}} \\
        \textbf{\gls{wd}} & Accuracy & Loss & Accuracy & Loss & Accuracy & F1 Score \\
        \hline
        1e-2 & 80.60 & 78.91 & 0.4387 & 0.4410 & 73.84 & 0.5950 \\
        1e-4 & 78.45 & 80.49 & 0.4390 & 0.4232 & 76.95 & 0.6740 \\
        1e-6 & 75.72 & 80.31 & 0.4688 & 0.4225 & 73.84 & 0.5747 \\
        0    & 78.59 & 80.24 & 0.4455 & 0.4151 & 77.72 & 0.6958 \\
        \hline
    \end{tabular}
    \caption{Average results for Resnet50 with \gls{rgb} images + \gls{ndvi} across different \gls{wd} values for a learning rate of 5e-6.}
    \label{tab:resnet50_wd_rgb_ndvi}
\end{table}

The results in Table \ref{tab:resnet50_wd_rgb_ndvi} indicate that weight decay does not 
significantly affect ResNet50's generalization when trained with \gls{rgb}+\gls{ndvi} 
images. While a \gls{wd} of $1\times10^{-2}$ achieves the highest validation accuracy 
(\textbf{80.60\%}), its test accuracy (\textbf{73.84\%}) is lower than the model trained 
without weight decay.

These findings suggest that \gls{wd} might not be the most effective regularization 
method for leveraging \gls{ndvi}-enhanced features for the Resnet50, as the best test results are obtained 
without any \gls{wd}.

\section{DenseNet161 Results}

Figure \ref{fig:d161_results} presents the training and validation performance for the input of
\gls{irg} images, showing trends in both loss and accuracy. Since the results for \gls{rgb} and \gls{rgb}+\gls{ndvi} 
images were very similar, only the \gls{irg} graphs are displayed for clarity.

\begin{figure}[h]
    \centering
    \includegraphics[width=1\textwidth]{../Figures/densenet161_averages.png}
    \caption{Average training and validation performance of Densenet161 across different \glspl{lr} with \gls{irg} images only.}
    \label{fig:d161_results}
\end{figure}

Compared to ResNet50's Figure \ref{fig:r50_results}, DenseNet shows more fluctuations 
in validation loss, even at lower \glspl{lr}. While higher \glspl{lr} 
still lead to instability, the expected smoothness at lower \glspl{lr} 
($1\times10^{-5}$ and $5\times10^{-6}$) is not as evident. This suggests that 
DenseNet may require additional regularization techniques.

Despite these fluctuations, the model still shows a general trend of improving 
accuracy over time. However, the inconsistent validation performance highlights 
the need for further tuning.

\subsection{IRG Results}

%                                                                     %
% ------------------------- DENSENET161 IRG ------------------------- %
%                                                                     %

The results for DenseNet161 trained on \gls{irg} images (Table \ref{tab:densenet161_irg}) show 
that lower \glspl{lr} contribute to improved validation performance 
and generalization. The best validation accuracy is achieved at a learning rate of 
$1\times10^{-5}$, reaching \textbf{82.76\%}, while the corresponding with a lower test accuracy 
of \textbf{77.72\%} with an F1 score of \textbf{0.7514}. 

\begin{table}[h]
    \centering
    \renewcommand{\arraystretch}{1.3}
    \setlength{\tabcolsep}{8pt}
    \begin{tabular}{|c|cccccc|}
        \hline
        & \multicolumn{2}{c}{\textbf{Validation}} & \multicolumn{2}{c}{\textbf{Train}} & \multicolumn{2}{c|}{\textbf{Test}} \\
        \textbf{\gls{lr}} & Accuracy & Loss & Accuracy & Loss & Accuracy & F1 Score \\
        \hline
        1e-4  & 78.74           & 0.4428 & 82.36 & 0.3940 & 72.80 & 0.7137 \\
        5e-5  & 80.03           & 0.4212 & 81.21 & 0.4149 & 76.17 & 0.6733 \\
        1e-5  & \textbf{82.76}  & 0.3994 & 83.87 & 0.3578 & 77.72 & 0.7514 \\
        5e-6  & 79.74           & 0.4454 & 82.18 & 0.3889 & 78.76 & 0.7617 \\
        \hline
    \end{tabular}
    \caption{Average results for Densenet161 with \gls{rgb} images across different \glspl{lr}.}
    \label{tab:densenet161_irg}
\end{table}

At an even smaller learning rate of $5\times10^{-6}$, the test accuracy reaches 
its highest value of \textbf{78.76\%}, along with a slightly improved F1 score of 
\textbf{0.7617}. This suggests that for this dataset and architecture, smaller learning rate allows 
for a more gradual and stable optimization process. 

However, higher \glspl{lr} such as $1\times10^{-4}$ and $5\times10^{-5}$ result 
in lower validation accuracy (\textbf{78.74\%} and \textbf{80.03\%}, respectively) and 
test accuracy (\textbf{72.80\%} and \textbf{76.17\%}), indicating that faster convergence 
may lead to suboptimal solutions.

Overall, the results indicate that \glspl{lr} in the range of $1\times10^{-5}$ to $5\times10^{-6}$ 
possess the best balance between training stability and generalization, producing the highest 
test performances.

\begin{table}[h]
    \centering
    \renewcommand{\arraystretch}{1.3}
    \setlength{\tabcolsep}{8pt}
    \begin{tabular}{|c|cccccc|}
        \hline
        & \multicolumn{2}{c}{\textbf{Validation}} & \multicolumn{2}{c}{\textbf{Train}} & \multicolumn{2}{c|}{\textbf{Test}} \\
        \textbf{\gls{wd}} & Accuracy & Loss & Accuracy & Loss & Accuracy & F1 Score \\
        \hline
        1e-2  & 81.03           & 0.4195 & 79.34 & 0.4387 & 82.47 & 0.7956 \\
        1e-4  & 80.17           & 0.4319 & 80.49 & 0.4193 & 81.35 & 0.7907 \\
        1e-6  & 77.59           & 0.4281 & 77.83 & 0.4558 & 78.76 & 0.7421 \\
        0     & \textbf{82.76}  & 0.3994 & 83.87 & 0.3578 & 77.72 & 0.7514 \\
        \hline
    \end{tabular}
    \caption{Average results for Densenet161 with \gls{irg} images across different \gls{wd} values for a learning rate of 5e-6.}
    \label{tab:densenet161_wd_irg}
\end{table}

The effect of \gls{wd} on DenseNet161's performance with \gls{irg} images is summarized in 
Table \ref{tab:densenet161_wd_irg}. Unlike the learning rate analysis, where lower values 
provided clear benefits, \gls{wd} shows a more complex impact on model performance.

A \gls{wd} of $1\times10^{-2}$ achieves the highest test accuracy (\textbf{82.47\%}) and the best 
test F1 score (\textbf{0.7956}), representing a significant improvement over the model trained 
without \gls{wd}. Interestingly, training without \gls{wd} yielded the highest validation 
accuracy (\textbf{82.76\%}).

At smaller \gls{wd} values, such as $1\times10^{-6}$, both validation accuracy
(\textbf{77.59\%}) and test accuracy (\textbf{78.76\%}) decrease, suggesting that insufficient 
regularization can reduce training stability. The unusually high validation accuracy observed 
in the no \gls{wd} setting may be an anomaly, despite being the average of multiple runs.


%                                                                     %
% ------------------------- DENSENET161 RGB ------------------------- %
%                                                                     %

\subsection{\gls{rgb} Results}

The results for DenseNet161 trained on \gls{rgb} images (Table \ref{tab:densenet161_rgb}) 
show once again that lower \glspl{lr} contribute to better generalization, following a 
similar trend observed in other configurations. The highest validation accuracy 
(\textbf{80.17\%}) is achieved at a learning rate of $5\times10^{-6}$, which also 
results in the best test accuracy (\textbf{78.76\%}) and a solid F1 score of 
\textbf{0.7485}.

\begin{table}[h]
    \centering
    \renewcommand{\arraystretch}{1.3}
    \setlength{\tabcolsep}{8pt}
    \begin{tabular}{|c|cccccc|}
        \hline
        & \multicolumn{2}{c}{\textbf{Validation}} & \multicolumn{2}{c}{\textbf{Train}} & \multicolumn{2}{c|}{\textbf{Test}} \\
        \textbf{\gls{lr}} & Accuracy & Loss & Accuracy & Loss & Accuracy & F1 Score \\
        \hline
        1e-4 & 78.59          & 0.4244 & 80.02 & 0.4248 & 76.94 & 0.7138 \\
        5e-5 & 76.58          & 0.4461 & 72.53  & 0.5401 & 74.36 & 0.7258 \\
        1e-5 & 77.44          & 0.4523 & 79.55  & 0.4285 & 75.13 & 0.7273 \\
        5e-6 & \textbf{80.17} & 0.4245 & 80.56 & 0.4098 & \textbf{78.76} & 0.7485 \\
        \hline
    \end{tabular}
    \caption{Average results for Densenet161 with \gls{rgb} images across different \glspl{lr}.}
    \label{tab:densenet161_rgb}
\end{table}

Higher \glspl{lr}, such as $1\times10^{-4}$ and $5\times10^{-5}$, yield lower 
validation accuracies (\textbf{78.59\%} and \textbf{76.58\%}, respectively) and 
suboptimal test performance (\textbf{76.94\%} and \textbf{74.36\%}). The instability 
in these models is further reflected in the higher validation loss values. 

The learning rate of $1\times10^{-5}$ achieves relatively balanced performance, with a 
validation accuracy of \textbf{77.44\%} and test accuracy of \textbf{75.13\%}. However, 
it does not outperform the lower learning rate of $5\times10^{-6}$. The results 
reinforce the idea that DenseNet161 benefits from a more controlled learning process, with 
$5\times10^{-6}$ being the most effective learning rate for \gls{rgb} inputs.


The effect of \gls{wd} on DenseNet161's performance with \gls{rgb} images is summarized 
in Table \ref{tab:densenet161_wd_rgb}. Unlike learning rate adjustments, where lower 
values clearly yielded better results, \gls{wd} presents a worse impact 
on generalization.

\begin{table}[h]
    \centering
    \renewcommand{\arraystretch}{1.3}
    \setlength{\tabcolsep}{8pt}
    \begin{tabular}{|c|cccccc|}
        \hline
        & \multicolumn{2}{c}{\textbf{Validation}} & \multicolumn{2}{c}{\textbf{Train}} & \multicolumn{2}{c|}{\textbf{Test}} \\
        \textbf{\gls{wd}} & Accuracy & Loss & Accuracy & Loss & Accuracy & F1 Score \\
        \hline
        1e-2     & 79.02   & 0.4195 & 79.12 & 0.4436  & 77.20  & 0.7442 \\
        1e-4   & 77.30     & 0.4685 & 80.27 & 0.4231  & 76.68  & 0.7239 \\
        1e-6 & 80.17       & 0.4049 & 78.76 & 0.4402  & 75.65  & 0.7403 \\
        0 & \textbf{80.17} & 0.4245 & 80.56 & 0.4098 & \textbf{78.76} & 0.7485 \\
        \hline
    \end{tabular}
    \caption{Average results for Densenet161 with \gls{rgb} images across different \gls{wd} values for a learning rate of 5e-6.}
    \label{tab:densenet161_wd_rgb}
\end{table}

Interestingly, training without \gls{wd} ($\gls{wd}=0$) yields the highest validation 
accuracy (\textbf{80.17\%}), the highest test accuracy (\textbf{78.76\%}), and 
a strong F1 score (\textbf{0.7485}). This suggests that for this particular dataset 
and model configuration, the \gls{wd} regularization may not be the best and could even 
slightly hinder performance.

A \gls{wd} of $1\times10^{-2}$ produces a slightly lower validation accuracy (\textbf{79.02\%}) 
and test accuracy (\textbf{77.20\%}), though its F1 score remains competitive at 
\textbf{0.7442}. This suggests that some more level of regularization can still 
provide stable performance without drastically reducing accuracy.

However, lower \gls{wd} values, such as $1\times10^{-6}$, do not necessarily 
lead to better generalization. While it achieves a decent validation accuracy 
(\textbf{80.17\%}), its test accuracy drops to \textbf{75.65\%}.


%                                                                            %
% ------------------------- DENSENET161 RGB + NDVI ------------------------- %
%                                                                            %

\subsection{\gls{rgb} + \gls{ndvi} Results}

The results for DenseNet161 trained on \gls{rgb} + \gls{ndvi} images (Table \ref{tab:densenet161_rgb_ndvi}) 
show the learning rate impact in the model's performance. The highest 
validation accuracy (\textbf{80.75\%}) is achieved at a learning rate of $1\times10^{-5}$, 
but this does not translate into the best test accuracy, which drops to \textbf{74.09\%} 
with an F1 score of \textbf{0.7195}.

\begin{table}[h]
    \centering
    \renewcommand{\arraystretch}{1.3}
    \setlength{\tabcolsep}{8pt}
    \begin{tabular}{|c|cccccc|}
        \hline
        & \multicolumn{2}{c}{\textbf{Validation}} & \multicolumn{2}{c}{\textbf{Train}} & \multicolumn{2}{c|}{\textbf{Test}} \\
        \textbf{\gls{lr}} & Accuracy & Loss & Accuracy & Loss & Accuracy & F1 Score \\
        \hline
        1e-4  & 78.88 & 0.4569 & 82.58 & 0.3855 & 78.24 & 0.7341 \\
        5e-5  & 79.31 & 0.4607 & 83.19 & 0.3710 & 79.28 & 0.7221 \\
        1e-5  & 80.75 & 0.4213 & 79.48 & 0.4377 & 74.09 & 0.7195 \\
        5e-6  & 79.60 & 0.4509 & 80.17 & 0.4166 & 77.72 & 0.7359 \\
        \hline
    \end{tabular}
    \caption{Average results for Densenet161 with \gls{rgb} images + \gls{ndvi} across different \glspl{lr}.}
    \label{tab:densenet161_rgb_ndvi}
\end{table}

The highest test accuracy (\textbf{79.28\%}) is achieved at a learning rate of $5\times10^{-5}$, 
which also yields a competitive validation accuracy of \textbf{79.31\%}. While a higher learning 
rate appears to improve test performance, the similarity in validation accuracies suggests that the 
improvement in test accuracy may be influenced by training dynamics rather than a consistently 
better generalization ability. This indicates potential sensitivity to hyperparameter selection 
or optimization behavior rather than a definitive advantage of a higher learning rate.


The impact of \gls{wd} on DenseNet161 trained with \gls{rgb} + \gls{ndvi} images is summarized 
in Table \ref{tab:densenet161_wd_rgb_ndvi}.

\begin{table}[h]
    \centering
    \renewcommand{\arraystretch}{1.3}
    \setlength{\tabcolsep}{8pt}
    \begin{tabular}{|c|cccccc|}
        \hline
        & \multicolumn{2}{c}{\textbf{Validation}} & \multicolumn{2}{c}{\textbf{Train}} & \multicolumn{2}{c|}{\textbf{Test}} \\
        \textbf{\gls{wd}} & Accuracy & Loss & Accuracy & Loss & Accuracy & F1 Score \\
        \hline
        1e-2  & 79.02  & 0.4358 & 79.12  & 0.4359  & 72.54  & 0.7006 \\
        1e-4  & 79.89  & 0.3906 & 80.49  & 0.4158  & 78.76  & 0.7515 \\
        1e-6  & 80.17  & 0.4287 & 79.91  & 0.4251  & 73.06  & 0.6977 \\
        0     & 79.60  & 0.4509 & 80.17  & 0.4166  & 77.72  & 0.7359 \\
        \hline
    \end{tabular}
    \caption{Average results for Resnet50 with \gls{rgb} images + \gls{ndvi} across different \gls{wd} values for a learning rate of 5e-6.}
    \label{tab:densenet161_wd_rgb_ndvi}
\end{table}

The results indicate that \gls{wd} 
does not have a strong or consistent effect on performance, as validation accuracy 
remains relatively stable across different values. While some variations exist in test 
accuracy and F1 score, the differences are minor and do not suggest a clear trend. This 
could indicate that the observed changes are primarily due to randomness rather than a 
meaningful influence of \gls{wd} on generalization.

\subsection{Gradual Unfreezing Strategy Results}

Given that the results obtained from different \glspl{lr} and \gls{wd} settings 
did not yield consistently strong performance, we experimented with a gradual unfreezing strategy, 
as described in the model architecture section, to explore potential improvements. The goal was 
to allow the model to retain useful pre-trained features while progressively fine-tuning deeper 
layers, thereby mitigating both underfitting and overfitting. This adjustment aimed to prevent 
excessive weight updates that could lead to the loss of valuable features learned in the initial layers.

The impact of this approach on DenseNet161 performance is presented in Table \ref{tab:best_runs}.
Compared to the previous results in Tables \ref{tab:densenet161_irg}, \ref{tab:densenet161_wd_irg}, 
\ref{tab:densenet161_rgb}, \ref{tab:densenet161_wd_rgb}, \ref{tab:densenet161_rgb_ndvi}
and \ref{tab:densenet161_wd_rgb_ndvi} unfreezing substantially improves both validation 
and test performance across all input types.

\begin{table*}[h]
    \centering
    \renewcommand{\arraystretch}{1.3}
    \setlength{\tabcolsep}{6pt}
    \begin{tabular}{|c|c|c|cc|cc|cc|}
        \hline
        \textbf{IMG} & \textbf{\gls{lr}} & \textbf{\gls{wd}} & \multicolumn{2}{c|}{\textbf{Validation}} & \multicolumn{2}{c|}{\textbf{Train}} & \multicolumn{2}{c|}{\textbf{Test}} \\
        &  &  & Accuracy & Loss & Accuracy & Loss & Accuracy & F1 Score \\
        \hline
        & 1e-4 & 1e-2 & 79.60 & 0.450 & 78.69 & 0.450 & 78.76 & 0.7485 \\
        \gls{rgb} & 1e-4 & 1e-4 & \textbf{83.05} & 0.386 & 78.40 & 0.478 & \textbf{84.46} & 0.8148 \\
        & 1e-5 & 1e-2 & 77.87 & 0.450 & 78.98 & 0.456 & 77.20 & 0.7349 \\
        & 1e-5 & 1e-4 & 79.60 & 0.438 & 81.57 & 0.411 & 76.68 & 0.7368 \\
        \hline
        & 1e-4 & 1e-2 & 78.74 & 0.431 & 78.62 & 0.454 & 80.31 & 0.7683 \\
        \gls{irg} & 1e-4 & 1e-4 & \textbf{80.46} & 0.426 & 82.65 & 0.396 & \textbf{82.90} & 0.7692 \\
        & 1e-5 & 1e-2 & 77.87 & 0.446 & 79.41 & 0.426 & 81.35 & 0.7778 \\
        & 1e-5 & 1e-4 & 76.72 & 0.498 & 78.76 & 0.450 & 75.65 & 0.7283 \\
        \hline
        & 1e-4 & 1e-2 & 81.90 & 0.387 & 79.27 & 0.438 & 77.20 & 0.7143 \\
        \gls{rgb}+\gls{ndvi} & 1e-4 & 1e-4 & \textbf{85.06} & 0.384 & 82.07 & 0.381 & \textbf{82.38} & 0.8000 \\ 
        & 1e-5 & 1e-2 & 82.76 & 0.405 & 79.70 & 0.435 & 77.72 & 0.7114 \\
        & 1e-5 & 1e-4 & 81.61 & 0.418 & 81.21 & 0.412 & 73.06 & 0.7079 \\
        \hline
    \end{tabular}
    \caption{Results for each \gls{lr} and \gls{wd} combination with gradual unfreezing strategy for DenseNet161.}
    \label{tab:best_runs}
\end{table*}



Starting with the \gls{rgb} setting, we notice the biggest improvement among them all, where test 
accuracy rises to \textbf{84.46\%} from the best previously achieved \textbf{78.76\%} 
(Table \ref{tab:densenet161_rgb}), an increase of over 5.7 percentage points. 
Similarly, the F1 score jumps from \textbf{0.7485} to \textbf{0.8148}, indicating 
not only better classification accuracy but also more balanced predictions.

For the \gls{irg} images, it shows a stable generalization improvement. Test accuracy 
reaches \textbf{82.90\%}, compared to a previous best of \textbf{82.47\%}
(Table \ref{tab:densenet161_wd_irg}). The associated F1-score is \textbf{0.7692}. 
Interestingly, another run achieved a slightly better F1-score of \textbf{0.7778} but at 
a lower test accuracy of \textbf{81.35\%}, indicating a trade-off between overall 
accuracy and class balance.

A similar trend is seen with \gls{rgb}+\gls{ndvi} images. The highest validation accuracy 
of \textbf{85.06\%} is now achieved, a notable increase from \textbf{80.75\%} in 
Table \ref{tab:densenet161_rgb_ndvi}. The test accuracy also improves
reaching \textbf{82.38\%}, which is much better than its best test 
accuracy before the unfreezing strategy (\textbf{79.28\%}).

Overall, these results demonstrate that gradually unfreezing is a highly effective 
strategy for DenseNet161, when applied to all types of imagery, it leads to 
significant performance gains in both validation and test accuracy.

\section{EfficientNet-B5 Results}


Figure \ref{fig:eb5_results} presents the training and validation performance for 
\gls{irg} images, showing trends in both loss and accuracy. Since the results for \gls{rgb} and \gls{rgb}+\gls{ndvi} 
images were similar, only the \gls{irg} graphs are displayed for clarity.

\begin{figure}[h]
    \centering
    \includegraphics[width=0.9\textwidth]{../Figures/efficientnet_b5_averages.png}
    \caption{Average training and validation performance of EfficientNetB5 across different \glspl{lr}.}
    \label{fig:eb5_results}
\end{figure}

The training loss curves indicate that models 
with higher \glspl{lr} ($\gls{lr}=1\times10^{-4}$, $\gls{lr}=5\times10^{-5}$) converge 
faster, whereas lower \glspl{lr} exhibit slower but steady progress. However, 
the validation loss plot shows that higher \glspl{lr} result in more volatility, 
suggesting that while they accelerate training, they may lead to instability in 
generalization, similarly to Resnet50.

%                                                                         %
% ------------------------- EFFICIENTNET_B5 IRG ------------------------- %
%                                                                         %
\subsection{\gls{irg} Results}

Table \ref{tab:efficientnet_b5_irg} presents results across different \glspl{lr}. 
While the highest validation accuracy (\textbf{79.74\%}) and test accuracy (\textbf{77.72\%}) were 
obtained with the largest \gls{lr} ($1\times10^{-4}$), the differences across learning 
rates remain relatively small. The best-performing configuration also achieved the highest training 
accuracy (\textbf{84.41\%}), suggesting strong learning on the training data. However, the 
validation loss (0.4542) and test F1-score (0.7584) indicating some signs of overfitting.

\begin{table}[h]
    \centering
    \renewcommand{\arraystretch}{1.3}
    \setlength{\tabcolsep}{8pt}
    \begin{tabular}{|c|cccccc|}
        \hline
        & \multicolumn{2}{c}{\textbf{Validation}} & \multicolumn{2}{c}{\textbf{Train}} & \multicolumn{2}{c|}{\textbf{Test}} \\
        \textbf{\gls{lr}} & Accuracy & Loss & Accuracy & Loss & Accuracy & F1 Score \\
        \hline
        1e-4  & 79.74  & 0.4542 & 84.41 & 0.3510 & 77.72 & 0.7584 \\
        5e-5  & 74.71  & 0.5184 & 77.50 & 0.4573 & 76.43 & 0.7425 \\
        1e-5  & 79.02  & 0.4590 & 78.47 & 0.4591 & 75.64 & 0.7054 \\
        5e-6  & 78.74  & 0.4473 & 79.01 & 0.4277 & 77.72 & 0.7415 \\
        \hline
    \end{tabular}
    \caption{Average results for EfficientNetB5 with \gls{irg} images across different \glspl{lr}.}
    \label{tab:efficientnet_b5_irg}
\end{table}

Given the close performance across \glspl{lr}, these results suggest that EfficientNet-B5 
may still require additional regularization to further stabilize validation and test performance. 
The minor variations in accuracy and F1-score highlight the potential influence of randomness, and 
tuning other hyperparameters, such as \gls{wd}, may offer better control over generalization.

Table \ref{tab:efficientnetb5_wd_irg} presents results across different \gls{wd} values
for the \gls{irg} images. The strongest test accuracy (\textbf{83.42\%}) and highest F1-score 
(\textbf{0.8182}) were obtained with $\gls{wd}=1\times10^{-2}$, demonstrating that stronger 
regularization improve generalization.

\begin{table}[h]
    \centering
    \renewcommand{\arraystretch}{1.3}
    \setlength{\tabcolsep}{8pt}
    \begin{tabular}{|c|cccccc|}
        \hline
        & \multicolumn{2}{c}{\textbf{Validation}} & \multicolumn{2}{c}{\textbf{Train}} & \multicolumn{2}{c|}{\textbf{Test}} \\
        \textbf{\gls{wd}} & Accuracy & Loss & Accuracy & Loss & Accuracy & F1 Score \\
        \hline
        1e-2 & 79.31    & 0.4374 & 80.42  & 0.4190  & \textbf{83.42}  & \textbf{0.8182} \\
        1e-4 & 77.59    & 0.4906 & 83.15  & 0.3831  & 80.83  & 0.7550 \\
        1e-6 & 79.60    & 0.4297 & 80.78  & 0.4086  & 80.31  & 0.7889 \\
        0    & 78.74    & 0.4473 & 79.01 & 0.4277 & 77.72 & 0.7415 \\
        \hline
    \end{tabular}
    \caption{Average results for EfficientNetB5 with \gls{irg} images across different \gls{wd} values for a learning rate of 5e-6.}
    \label{tab:efficientnetb5_wd_irg}
\end{table}

Other \gls{wd} values exhibit varying degrees of effectiveness. A smaller \gls{wd} 
($\gls{wd}=1\times10^{-4}$) led to high training accuracy of \textbf{83.15\%} but a comparatively lower 
test accuracy \textbf{80.83\%}. Meanwhile, an extremely small \gls{wd} ($\gls{wd}=1\times10^{-6}$) 
still achieved strong generalization with a test accuracy of \textbf{80.31\%} and an F1-score of 0.7889.

Overall, these results demonstrate that for EfficientNet-B5 with \gls{irg} images, 
regularization plays a crucial role in improving test accuracy and generalization. 
While learning rate adjustments alone had limited impact, a combination of a learning 
rate of $5\times10^{-6}$ and a \gls{wd} of $1\times10^{-2}$ produced very good 
results.


%                                                                         %
% ------------------------- EFFICIENTNET_B5 \gls{rgb} ------------------------- %
%                                                                         %

\subsection{\gls{rgb} Results}
The different learning rate results for EfficientNet-B5 trained on \gls{rgb} images are 
presented in Table \ref{tab:efficientnetb5_lr_rgb}. The best validation accuracy is 
achieved with $\gls{lr}=1\times10^{-5}$ at 78.45\%, while the highest test accuracy is 
obtained with the lowest learning rate ($\gls{lr}=5\times10^{-6}$), reaching 78.76\%. The best 
F1-score is found at $\gls{lr}=1\times10^{-4}$ with a value of 0.7616.

However, similar to the results obtained for \gls{irg} images (Table \ref{tab:efficientnet_b5_irg}),
there is no significant improvement across \glspl{lr}, and overall test performance 
remains suboptimal. The results suggest that while lower \glspl{lr} may provide slight 
improvements in generalization, the model still struggles to achieve strong validation and 
test performance.

\begin{table}[h]
    \centering
    \renewcommand{\arraystretch}{1.3}
    \setlength{\tabcolsep}{8pt}
    \begin{tabular}{|c|cccccc|}
        \hline
        & \multicolumn{2}{c}{\textbf{Validation}} & \multicolumn{2}{c}{\textbf{Train}} & \multicolumn{2}{c|}{\textbf{Test}} \\
        \textbf{\gls{lr}} & Accuracy & Loss & Accuracy & Loss & Accuracy & F1 Score \\
        \hline
        1e-4   & 76.29  & 0.4923 & 79.05 & 0.4659 & 78.24 & 0.7616 \\
        5e-5   & 78.16  & 0.4519 & 77.14 & 0.4834 & 77.72 & 0.7387 \\
        1e-5   & 78.45  & 0.4797 & 79.19 & 0.4467 & 74.10 & 0.7367 \\
        5e-6   & 75.86  & 0.4725 & 77.29 & 0.4635 & 78.76 & 0.7360 \\
        \hline
    \end{tabular}
    \caption{Average results for EfficientNetB5 with \gls{irg} images across different \glspl{lr}.}
    \label{tab:efficientnetb5_lr_rgb}
\end{table}

Given these observations, we next explore the impact of \gls{wd} on EfficientNet-B5's 
performance. Since learning rate adjustments alone did not lead to consistent improvements, 
applying different \gls{wd} values may help introduce better regularization and stability to the model
the same way it did for \gls{irg} images.

The Table \ref{tab:efficientnetb5_wd_rgb} presents the results for different 
\gls{wd} values. The highest validation accuracy is observed at 
$\gls{wd}=1\times10^{-4}$ with \textbf{81.90\%}. However, the best test accuracy is 
achieved at $\gls{wd}=1\times10^{-2}$, reaching \textbf{81.87\%}. The highest F1-score 
is also obtained at $\gls{wd}=1\times10^{-2}$ with 0.7619. These results indicate 
that stronger \gls{wd} improves test performance, even if it does not 
always yield the highest validation accuracy.

\begin{table}[h]
    \centering
    \renewcommand{\arraystretch}{1.3}
    \setlength{\tabcolsep}{8pt}
    \begin{tabular}{|c|cccccc|}
        \hline
        & \multicolumn{2}{c}{\textbf{Validation}} & \multicolumn{2}{c}{\textbf{Train}} & \multicolumn{2}{c|}{\textbf{Test}} \\
        \textbf{\gls{wd}} & Accuracy & Loss & Accuracy & Loss & Accuracy & F1 Score \\
        \hline
        1e-6  & 78.16           & 0.4763 & 76.39  & 0.4937  & 81.35  & 0.7534 \\
        0     & 75.86           & 0.4725 & 77.29  & 0.4635 & 78.76 & 0.7360 \\
        1e-2  & 81.63           & 0.4344 & 78.47  & 0.4420  & \textbf{81.87}  & 0.7619 \\
        1e-4  & \textbf{81.90}  & 0.4064 & 79.27  & 0.4505  & 80.83 & 0.7413 \\
        \hline
    \end{tabular}
    \caption{Average results for EfficientNetB5 with \gls{rgb} images across different \gls{wd} values for a learning rate of 5e-6.}
    \label{tab:efficientnetb5_wd_rgb}
\end{table}


%                                                                                %
% ------------------------- EFFICIENTNET_B5 \gls{rgb} + \gls{ndvi} ------------------------- %
%                                                                                %

\subsection{RGB + NDVI Results}

Table \ref{tab:efficientnetb5_lr_rgb_ndvi} presents the results for EfficientNet-B5 trained on 
\gls{rgb} images + \gls{ndvi} with different \glspl{lr}. The highest validation accuracy (\textbf{81.47\%}) is 
achieved at $\gls{lr}=5\times10^{-5}$, but this configuration also results in the lowest test 
accuracy (70.73\%). In contrast, the best test accuracy (\textbf{79.83\%}) and highest F1-score 
(0.7886) are obtained with the lowest learning rate, $\gls{lr}=5\times10^{-6}$.

However, similar to the previous results for \gls{rgb} and \gls{irg} images, the overall performance 
remains inconsistent across different \glspl{lr}, with no clear pattern indicating 
an optimal setting. The variability in validation and test accuracy suggests that 
regularization techniques, such as \gls{wd}, may be necessary to stabilize generalization.

\begin{table}[h]
    \centering
    \renewcommand{\arraystretch}{1.3}
    \setlength{\tabcolsep}{8pt}
    \begin{tabular}{|c|cccccc|}
        \hline
        & \multicolumn{2}{c}{\textbf{Validation}} & \multicolumn{2}{c}{\textbf{Train}} & \multicolumn{2}{c|}{\textbf{Test}} \\
        \textbf{\gls{lr}} & Accuracy & Loss & Accuracy & Loss & Accuracy & F1 Score \\
        \hline
        1e-4   & 78.02 &  0.4699 & 80.13 & 0.4456 & 75.39 & 0.7400 \\
        5e-5   & 81.47 &  0.4405 & 80.99 & 0.4185 & 70.73 & 0.7064 \\
        1e-5   & 77.87 &  0.4623 & 80.92 & 0.4181 & 73.05 & 0.7475 \\
        5e-6   & 75.29 &  0.4815 & 80.42 & 0.4354 & 79.83 & 0.7886 \\
        \hline
    \end{tabular}
    \caption{Average results for EfficientNet-B5 with \gls{rgb}+\gls{ndvi} images across different \glspl{lr}.}
    \label{tab:efficientnetb5_lr_rgb_ndvi}
\end{table}

To address the instability observed with different \glspl{lr}, we next evaluate the impact 
of \gls{wd}. Table \ref{tab:efficientnetb5_wd_rgb_ndvi} presents the results across different 
\gls{wd} values for a fixed learning rate of $5\times10^{-6}$. The highest validation 
accuracy (\textbf{80.17\%}) and test accuracy (\textbf{81.35\%}) are obtained with $\gls{wd}=1\times10^{-4}$, which also 
achieves the highest F1-score (\textbf{0.7955}). This suggests that moderate \gls{wd} helps improve 
generalization and stability.

\begin{table}[h]
    \centering
    \renewcommand{\arraystretch}{1.3}
    \setlength{\tabcolsep}{8pt}
    \begin{tabular}{|c|cccccc|}
        \hline
        & \multicolumn{2}{c}{\textbf{Validation}} & \multicolumn{2}{c}{\textbf{Train}} & \multicolumn{2}{c|}{\textbf{Test}} \\
        \textbf{\gls{wd}} & Accuracy & Loss & Accuracy & Loss & Accuracy & F1 Score \\
        \hline
        1e-2  & 76.72  & 0.4617  & 77.90  & 0.4819  & 74.09  & 0.6154 \\
        1e-4  & \textbf{80.17}  & 0.4309   & 80.42  & 0.4221  & \textbf{81.35}  & \textbf{0.7955} \\
        1e-6  & 75.57  & 0.4838  & 78.76  & 0.4466  & 79.79  & 0.7845 \\
        0     & 75.29 & 0.4815 & 80.42 & 0.4354 & 79.83 & 0.7886 \\
        \hline
    \end{tabular}
    \caption{Average results for EfficientNet-B5 with \gls{rgb}+\gls{ndvi} images across different \gls{wd} values for a learning rate of 5e-6.}
    \label{tab:efficientnetb5_wd_rgb_ndvi}
\end{table}

\section{Summary of Best Results for Each Architecture}

The performance of different architectures (ResNet50, DenseNet161, and EfficientNet-B5) 
across the three input configurations (\gls{rgb}, \gls{irg}, and \gls{rgb}+\gls{ndvi}) show some variations 
in accuracy and generalization capabilities. Table \ref{tab:model_comparison} summarizes the best 
results for each architecture.

\begin{table}[h]
    \centering
    \renewcommand{\arraystretch}{1.3}
    \setlength{\tabcolsep}{8pt}
    \begin{tabular}{|c|c|c|c|c|}
        \hline
        \textbf{Model} & \textbf{Image Type} & \textbf{Val Acc (\%)} & \textbf{Test Acc (\%)} & \textbf{Test F1 (\%)} \\
        \hline
                 & \gls{irg}        & \textbf{81.32}  & \textbf{82.90}  & \textbf{80.24} \\
        ResNet50 & \gls{rgb}        & 82.18  & 80.83  & 77.30 \\
                 & \gls{rgb}+\gls{ndvi}   & 78.59  & 77.72  & 69.58 \\
        \hline
                    & \gls{irg}        & 80.46  & 82.90  & 76.92 \\
        DenseNet161 & \gls{rgb}        & \textbf{83.05}  & \textbf{84.46}  & \textbf{81.48} \\
                    & \gls{rgb}+\gls{ndvi}   & 85.06  & 82.38  & 80.00 \\
        \hline
                        & \gls{irg}        & 79.31  & \textbf{83.42}  & \textbf{81.82} \\
        EfficientNet-B5 & \gls{rgb}        & 81.63  & 81.87  & 76.19 \\
                        & \gls{rgb}+\gls{ndvi}   & 80.17  & 81.35  & 79.55 \\
        \hline
    \end{tabular}
    \caption{Comparison of the best-performing configurations for ResNet50, DenseNet161, and EfficientNet-B5 across different image types.}
    \label{tab:model_comparison}
\end{table}


\subsection{ResNet50 Performance}
For ResNet50, the \gls{irg} configuration yielded the best overall performance, with the highest 
test accuracy (\textbf{82.90\%}) and F1 score (\textbf{80.24\%}). The \gls{rgb} configuration performed 
slightly worse, achieving a validation accuracy of \textbf{82.18\%}, but a lower test accuracy 
(\textbf{80.83\%}) and F1 score (\textbf{77.30\%}). \gls{rgb}+\gls{ndvi} produced the weakest results across all 
metrics, suggesting that adding \gls{ndvi} does not enhance ResNet50's ability to generalize.

The poor performance of the \gls{rgb}+\gls{ndvi} configuration may be attributed 
to the architectural design of ResNet50. Originally designed for three-channel 
\gls{rgb} images, the model was adapted for four-channel input, but this adaptation 
may not have been sufficient for effectively leveraging the \gls{ndvi} information. 

One possible reason for this lower performance is inefficient channel utilization, 
as the ResNet50 architecture is primarily designed for three-channel inputs. The addition 
of a fourth channel may not be fully leveraged by the initial convolutional layers.

\subsection{DenseNet161 Performance}
DenseNet161 showed strong generalization, with \gls{rgb} achieving the highest test accuracy (\textbf{84.46\%}) of all results
and \gls{rgb}+\gls{ndvi} attaining the best validation accuracy (\textbf{85.06\%}). However, despite excelling in validation, 
\gls{rgb}+\gls{ndvi} underperformed in test accuracy (\textbf{82.38\%}). The \gls{irg} configuration also remained competitive, 
reaching a test accuracy of \textbf{82.90\%}.

\subsection{EfficientNet-B5 Performance}
EfficientNet-B5 exhibited its strongest performance with \gls{irg}, achieving the highest test accuracy (\textbf{83.42\%}) 
and F1 score (\textbf{81.82\%}), indicating its effectiveness in capturing spectral features. While \gls{rgb} and \gls{rgb}+\gls{ndvi} 
configurations performed slightly worse (\textbf{81.87\%} and \textbf{81.35\%} test accuracy, respectively), 
their close results suggest EfficientNet-B5 maintains relatively stable performance across different input configurations.

\section{Ensemble Model Results}

In this section, we present the results of the ensemble models, which combine the predictions of
the best-performing models for each architecture and input type. Firstly we present the results of the
\gls{rgb}-Only Ensemble, which combines the best-performing models trained on \gls{rgb} images for each architecture.
Next, we present the results of the Ensemble that combines the best models for different input
types (\gls{rgb}, \gls{irg}, and \gls{rgb}+\gls{ndvi}).

\subsection{\gls{rgb} Only Ensemble}

The \gls{rgb}-Only Ensemble was evaluated based on accuracy and F1-score, achieving the results
presented in Table \ref{tab:rgb_ensemble_results}. The ensemble achieved an accuracy of \textbf{84.97\%} and an F1-score of \textbf{84.99\%}.

\begin{table}[h]
    \centering
    \renewcommand{\arraystretch}{1.3}
    \setlength{\tabcolsep}{8pt}
    \begin{tabular}{|c|c|}
        \hline
        \textbf{Metric} & \textbf{Ensemble Score} \\
        \hline
        Accuracy & 84.97\% \\
        F1-Score & 84.99\% \\
        \hline
    \end{tabular}
    \caption{Performance metrics of the \gls{rgb}-Only Ensemble.}
    \label{tab:rgb_ensemble_results}
\end{table}

To further evaluate the ensemble's classification performance, Figure \ref{fig:rgb_confusion} presents 
the confusion matrix. The normalized confusion matrix provides insight into the model's ability to 
correctly classify each class.

\begin{figure}[h]
    \centering
    \includegraphics[width=0.9\textwidth]{confusion_matrix_rgb_ensemble.png}
    \caption{Normalized Confusion Matrix for the \gls{rgb}-Only Ensemble.}
    \label{fig:rgb_confusion}
\end{figure}

The confusion matrix indicates that the ensemble correctly classified \textbf{86\% of Average Risk} cases 
and \textbf{83\% of Very High Risk} cases. The misclassification rates of 14\% and 17\%, respectively, 
suggest that while the ensemble achieves strong performance, improvements could be explored through 
alternative ensembling techniques or adjustments in decision thresholds.

\subsection{Multi-Input Ensemble}

The Multi-Input Ensemble combines models trained with different input configurations (\gls{rgb}, \gls{irg}, and \gls{rgb}+\gls{ndvi}) 
across DenseNet161 and EfficientNet-B5. This approach aims to leverage complementary spectral information to 
improve classification performance. 

ResNet50 with \gls{rgb}+\gls{ndvi} was excluded from this ensemble due to its lower test accuracy and F1-score compared 
to other models (Table \ref{tab:model_comparison}), ensuring that only the most effective configurations 
contributed to the final predictions.

\begin{table}[h]
    \centering
    \renewcommand{\arraystretch}{1.3}
    \setlength{\tabcolsep}{8pt}
    \begin{tabular}{|c|c|}
        \hline
        \textbf{Metric} & \textbf{Ensemble Score} \\
        \hline
        Accuracy & 86.53\% \\
        F1-Score & 86.53\% \\
        \hline
    \end{tabular}
    \caption{Performance metrics of the Multi-Input Ensemble.}
    \label{tab:multi_input_ensemble_results}
\end{table}

The Multi-Input Ensemble achieved a very strong F1-score of \textbf{86.53\%}, indicating a well-balanced trade-off between 
precision and recall. This result suggests that the ensemble is not only making accurate predictions but is also 
maintaining a consistent ability to correctly classify both risk categories. 

The high F1-score is particularly important in this context, as it reflects the model's ability to minimize both 
false positives (misclassifying Average Risk as Very High Risk) and false negatives (failing to identify Very High 
Risk cases).

Figure \ref{fig:multi_input_confusion} presents the normalized confusion matrix for the Multi-Input 
Ensemble. This matrix provides insight into classification performance across the two risk categories.

\begin{figure}[h]
    \centering
    \includegraphics[width=0.9\textwidth]{confusion_matrix_ensemble.png}
    \caption{Normalized Confusion Matrix for the Multi-Input Ensemble.}
    \label{fig:multi_input_confusion}
\end{figure}

The confusion matrix indicates that the ensemble correctly classified \textbf{88\% of Average Risk} cases 
and \textbf{84\% of Very High Risk} cases. The misclassification rates of 12\% and 16\%, respectively, 
affirm that this approach slightly outperforms the \gls{rgb}-Only Ensemble in distinguishing samples. 
This improvement demonstrates the benefit of integrating multiple input types, 
as different spectral channels likely contributed to better feature extraction and classification.

\subsection{Analysis of High-Confidence Misclassifications}

One critical issue is the set of misclassified samples where the ensemble model predicted with 
\textbf{100\% confidence} but failed to match the ground truth. These cases are concerning as they reveal 
potential weaknesses in the model's learned features or decision-making process.

To analyze these misclassifications, we present both the \gls{rgb} and \gls{irg} representations of the 
affected locations. Below, two houses are shown where the ensemble classified them as \textbf{Average Risk}, 
despite their ground truth label being \textbf{Very High Risk},a potentially dangerous error 
in real-world applications.


\begin{figure}[ht]
    \centering
    \begin{minipage}{0.75\textwidth}
        \centering
        \begin{minipage}[b]{0.49\textwidth}
            \centering
            \includegraphics[width=\textwidth]{../Figures/BE47.png}
            \caption{False Negative Example 1 - \gls{rgb}}
            \label{fig:fn1_rgb}
        \end{minipage}
        \hfill
        \begin{minipage}[b]{0.49\textwidth}
            \centering
            \includegraphics[width=\textwidth]{../Figures/BE47_IRG.png}
            \caption{False Negative Example 1 - \gls{irg}}
            \label{fig:fn1_irg}
        \end{minipage}

        \begin{minipage}[b]{0.49\textwidth}
            \centering
            \includegraphics[width=\textwidth]{../Figures/V29.png}
            \caption{False Negative Example 2 - \gls{rgb}}
            \label{fig:fn2_rgb}
        \end{minipage}
        \hfill
        \begin{minipage}[b]{0.49\textwidth}
            \centering
            \includegraphics[width=\textwidth]{../Figures/V29_IRG.png}
            \caption{False Negative Example 2 - \gls{irg}}
            \label{fig:fn2_irg}
        \end{minipage}
    \end{minipage}
    
    \caption{False Negative Cases: Houses classified as \textit{Average Risk} but labeled as \textit{Very High Risk}.}
    \label{fig:false_negatives}
\end{figure}


After closely reviewing these two samples, they show visual characteristics typically associated with 
lower-risk properties (Figure \ref{fig:false_negatives}).

The first example is a well-maintained house with minimal vegetation in its immediate 
surroundings. The area appears to be clean and managed, with little to no overgrown vegetation. 
However, the labeling may have been particularly strict in this case, as there is a tree and 
some vegetation close to the structure, which might have contributed to its classification as 
\textbf{Very High Risk}.

In the second example, while vegetation density is slightly higher, the house itself remains 
well-kept, and the surrounding vegetation appears to be managed rather than overgrown. Similar 
to the first case, the proximity of vegetation to the house may have influenced the labeling decision, 
leading to its classification as \textbf{Very High Risk}.


Below, we present two examples where houses were labeled as 
\textbf{Average Risk} but were predicted as \textbf{Very High Risk} by the model.

\begin{figure}[ht]
    \centering
    \begin{minipage}{0.75\textwidth}
        \centering
        \begin{minipage}[b]{0.49\textwidth}
            \centering
            \includegraphics[width=\textwidth]{../Figures/V6.png}
            \caption{False Positive Example 1 - \gls{rgb}}
            \label{fig:fp1_rgb}
        \end{minipage}
        \hfill
        \begin{minipage}[b]{0.49\textwidth}
            \centering
            \includegraphics[width=\textwidth]{../Figures/V6_IRG.png}
            \caption{False Positive Example 1 - \gls{irg}}
            \label{fig:fp1_irg}
        \end{minipage}

        \vspace{0.1cm}

        \begin{minipage}[b]{0.49\textwidth}
            \centering
            \includegraphics[width=\textwidth]{../Figures/VC135.png}
            \caption{False Positive Example 2 - \gls{rgb}}
            \label{fig:fp2_rgb}
        \end{minipage}
        \hfill
        \begin{minipage}[b]{0.49\textwidth}
            \centering
            \includegraphics[width=\textwidth]{../Figures/VC135_IRG.png}
            \caption{False Positive Example 2 - \gls{irg}}
            \label{fig:fp2_irg}
        \end{minipage}
    \end{minipage}
    
    \caption{False Positive Cases: Houses classified as \textit{Very High Risk} but labeled as \textit{Average Risk}.}
    \label{fig:false_positives}
\end{figure}


After reviewing these two cases (Figure \ref{fig:false_positives}), both exhibit visual characteristics 
strongly associated with high fire risk, despite being labeled as \textbf{Average Risk} in the dataset.

In the first example, the property is closely surrounded by dense vegetation, including trees and 
shrubs, which are key indicators of elevated fire risk. 

In the second example, unlike the first case where vegetation proximity is the primary concern, 
this location features widespread, uninterrupted vegetation cover. According to Table \ref{fig:fire-risk-table}, 
extensive vegetation continuity is a critical parameter for extreme fire risk, which likely contributed 
to the ensemble model's classification of Very High Risk.

Another observations is that shadows present in the imagery may have influenced the model's perception of 
vegetation density, potentially leading to an overestimation of fire risk. In areas where dense tree 
canopies cast dark shadows, the model might interpret these regions as continuous vegetation cover rather 
than open or managed spaces. This effect could have contributed to the ensemble's misclassification of certain 
as Very High Risk despite their actual fire exposure conditions.
